% =============================================================================
% Chapter 5: KV Cache Management — Agentic ML Systems Lecture Slides
% =============================================================================
% Built around the chapter thesis:
%   "Retained attention state, and how long it stays, sets agent serving
%    capacity."
% One running case (the repair agent's r18 request: 32,768 tokens, then a
% 45-second integration test, served by Llama 3 70B on one 8-accelerator
% node) carries the deck from bytes per token, through allocation, sharing,
% and reuse, to the keep/move/drop decision at the tool wait and the turn
% rate that residence time allows.
%
% Slides are written to be read without a speaker: each frame states its
% claim in the title, shows the evidence, and closes with a bottom line.
% Speaker notes add pacing, questions, and misconceptions.
%
% Numbers are the rendered chapter's computed values (LEGO cells backed by
% MLSysIM) at the commit this deck was built against.
% =============================================================================
\documentclass[aspectratio=169, 12pt]{beamer}
\usepackage{../../assets/beamerthememlsys}

\mlsyssetup{
  volume       = {Volume III},
  chapter      = {Chapter 5},
  logo         = {../../assets/img/logo-mlsysbook.png},
  instlogo     = {../../assets/img/logo-harvard.png},
  chaptertitle = {KV Cache Management},
}

% --- Fonts ---
\usepackage[T1]{fontenc}
\usepackage[scaled=0.9]{helvet}
\usepackage{courier}
\renewcommand{\familydefault}{\sfdefault}

% --- Packages ---
\usepackage{booktabs}
\usepackage{amsmath}
\usepackage{array}
\newcolumntype{L}[1]{>{\raggedright\arraybackslash}p{#1}}

% --- Image paths ---
\graphicspath{
  {images/}
}

% --- Chapter-specific macros (glyphs follow the book's notation registry) ---
\newcommand{\mtok}{m_{\text{token}}}
\newcommand{\Mkv}{M_{\text{KV}}}
\newcommand{\Mpool}{M_{\text{pool}}}
\newcommand{\NL}{N_L}
\newcommand{\Hkv}{H_{\text{KV}}}
\newcommand{\dhead}{d_{\text{head}}}
\newcommand{\selem}{s_{\text{elem}}}
\newcommand{\Sblk}{S_{\text{blk}}}
\newcommand{\Spre}{S_{\text{prefix}}}
\newcommand{\Kmax}{K_{\max}}
\newcommand{\fhead}{f_{\text{head}}}
\newcommand{\BWlink}{\text{BW}_{\text{link}}}
\newcommand{\Rattain}{R_{\text{attain}}}
\newcommand{\Ttool}{T_{\text{tool}}}
\newcommand{\Tmodel}{T_{\text{model}}}
\newcommand{\Nslot}{\bar{N}_{\text{slot}}}
\newcommand{\lam}{\lambda_{\text{arr}}}
\newcommand{\Thold}{\bar{T}_{\text{hold}}}
\newcommand{\Nbudget}{N_{\text{budget}}}
% Bottom line: the one sentence a reader without a speaker should keep.
\newcommand{\bottomline}[1]{%
  \vfill
  \begin{mlsyscard}[top=0.15em,bottom=0.15em,fontupper=\footnotesize]{crimson}\textbf{Bottom line.} #1\end{mlsyscard}%
}
\definecolor{bookblue}{HTML}{006395}
\definecolor{bookorange}{HTML}{E67817}
\definecolor{bookgreen}{HTML}{008F45}
\definecolor{bookviolet}{HTML}{7E317B}
\definecolor{bookred}{HTML}{CB202D}

% Compact reference entry for a long reference list
\newcommand{\mlsysrefsmall}[2]{{\footnotesize\textcolor{crimson}{\textbf{[#1]}} #2}\\[0.01cm]}

% --- Helper: safe image include ---
\newcommand{\safeimg}[2][width=\textwidth,keepaspectratio]{%
  \IfFileExists{images/#2}{\includegraphics[#1]{#2}}{%
    \IfFileExists{#2}{\includegraphics[#1]{#2}}{%
      \fbox{\parbox[c][2.5cm][c]{0.85\linewidth}{%
        \centering\footnotesize\textcolor{midgray}{[Missing image]}}}%
    }%
  }%
}

% --- Section count (must match actual \section{} count) ---
\setcounter{mlsystotalsections}{8}

\title{KV Cache Management}
\author{Vijay Janapa Reddi}
\institute{Harvard University}
\date{}

\begin{document}

% =============================================================================
% TITLE SLIDE
% =============================================================================
\mlsystitle{KV Cache Management}{Retained attention state, and how long\\it stays, sets agent serving capacity.}{cover_virtual_memory_blueprint.png}

% =============================================================================
% LEARNING OBJECTIVES
% =============================================================================
\begin{frame}{Learning Objectives}
\note{
[1 min]
% -- NARRATE
Read the subtitle aloud. It has two halves, and the lecture follows them:
first how many bytes the state takes and how to pack them (objectives 1 to
5), then how long each byte stays and what that does to the turn rate
(objectives 6 and 7).
% -- FLEX
[CORE] Do not dwell; return to this list at the wrap-up.
}

\small
By the end of this lecture you should be able to:
\begin{enumerate}\setlength{\itemsep}{0pt}
  \item \textbf{Calculate} KV memory and \textbf{resident fit} from model geometry, tokens, and precision.
  \item \textbf{Calculate} the memory \textbf{paged blocks} and \textbf{prefix sharing} recover.
  \item \textbf{Distinguish} compatible, available, permitted reuse from \emph{current} evidence.
  \item \textbf{Calculate} the token-weighted \textbf{prefix hit rate} of a radix prefix cache.
  \item \textbf{Select} a \textbf{chunked-prefill} budget that bounds the gap between output tokens.
  \item \textbf{Evaluate} keep, move, drop, and lossy approximation for paused state.
  \item \textbf{Analyze} how residence time sets turn rate: \textbf{Little's law}, the \textbf{residency tax}.
\end{enumerate}
\end{frame}

% =============================================================================
% SECTION 1: FOOTPRINT
% =============================================================================
\section{Footprint}

\begin{frame}{An agent generating nothing can still fill the node}
\note{
[3 min]
% -- LINK
The tool-wait problem was named in the turn lifecycle at the start of the
book. Today it gets a byte count and a price.
% -- NARRATE
Tell the r18 story: the repair agent's reproducer passes, the harness sends
a 32,768-token request, then launches a 45-second integration test. During
the test the model produces nothing. Ask where its attention state is. Then
point to the right card: the caller and the serving operator hold different
levers, and with a hosted API the caller never touches accelerator memory.
% -- ENGAGE
Ask: ``During those 45 seconds, what is the accelerator doing for this
request?'' [Nothing, yet its memory is still claimed.]
% -- WARN
Students hear ``runtime'' and think of the serving engine. In this volume
the runtime is the agent runtime around the loop; the inference software is
the serving engine.
% -- FLEX
[CORE]
}

\begin{columns}[T]
\begin{column}{0.53\textwidth}
\small
\textbf{Running case: the repair agent at r18}
\begin{itemize}\setlength{\itemsep}{1pt}
  \item the harness sends a \textbf{32,768-token} request
  \item then runs an integration test for \textbf{45 s}
  \item the model generates \textbf{zero} tokens meanwhile
  \item its attention state can stay resident throughout
\end{itemize}
\vspace{0.05cm}
The serving engine must keep the state, move it to slower memory, or drop it and rebuild it.
\end{column}
\begin{column}{0.43\textwidth}
\begin{mlsyscard}{bookviolet}
\footnotesize\textbf{Two parties, two sets of levers}\\[2pt]
\textbf{Caller (harness):} context, tool timing, any cache hints the API exposes.\\[2pt]
\textbf{Serving operator:} allocation, placement, scheduling, retention.
\end{mlsyscard}
\end{column}
\end{columns}

\bottomline{Retained attention state is a claim on accelerator memory that persists while the agent waits, so how long it stays matters as much as how big it is.}
\end{frame}

\begin{frame}{Each retained token stores a key and a value at every layer}
\note{
[3 min]
% -- LINK
Autoregressive decoding and its retained state were introduced with the
foundation model; here we only count the bytes.
% -- NARRATE
Walk the bar: system prompt, tool schemas, workspace, turn history, latest
observation. Then say that none of this structure survives: every token
becomes a key and a value vector at every one of the 80 layers, for each of
8 key-value heads. Do the arithmetic: 2 times 80 times 8 times 128 times 2
bytes is 327,680 bytes per token; times 32,768 tokens is 10.7 GB.
% -- ENGAGE
Ask: ``Prefill finishes. Is this request at its peak memory?'' [No. Every
output token appends a key and value at every layer, so the peak is at
completion.]
% -- WARN
Students budget for the prompt only. The definition's pitfall is exactly
that: output tokens grow the cache too.
% -- FLEX
[CORE]
}

\centering
\safeimg[width=0.95\textwidth,height=2.35cm,keepaspectratio]{context-zones.pdf}

\vspace{0.05cm}
{\large $\mtok = 2 \cdot \NL \cdot \Hkv \cdot \dhead \cdot \selem \qquad \Mkv = S \cdot \mtok$}

\vspace{0.1cm}
\raggedright\small
\textbf{Llama 3 70B, 16-bit:} $2 \times 80 \times 8 \times 128 \times 2 = \mathbf{327{,}680}$ bytes per token; for 32,768 tokens, $\Mkv \approx \mathbf{10.7}$ \textbf{GB}. The \textbf{KV cache} keeps these keys and values so later tokens need not recompute them.

\bottomline{The engine sees bytes, not zones: r18's 32,768 tokens cost 10.7 GB whatever they say.}
\end{frame}

\begin{frame}{The pool left after the weights holds 48 r18 contexts}
\note{
[3 min]
% -- NARRATE
Build the pool on the board: 8 accelerators of 80 GB, which is 85.9 GB in
decimal units, give 687.2 GB. Weights take 141.2 GB, the runtime reserve 20
GB, leaving 526.0 GB. Divide by 10.74 GB and take the floor: 48. Then read
the curve: every fourfold increase in context cuts the count by about three
quarters.
% -- ENGAGE
Ask: ``Who controls where on this curve the service sits?'' [The agent
runtime, through how many tokens it retains.]
% -- WARN
These are raw tensor payloads, balanced across shards. Allocator overhead,
fragmentation, and headroom lower them, and every single device must fit
its shard.
% -- FLEX
[CORE]
}

\begin{columns}[T]
\begin{column}{0.47\textwidth}
\small
\textbf{The KV pool} (equation 5)\\[2pt]
$\Mpool = 687.2 - 141.2 - 20 \approx 526.0$ GB\\
\footnotesize(8 accelerators, weights, runtime reserve)

\vspace{0.15cm}
\small\textbf{Resident fit} $= \lfloor \Mpool / \Mkv \rfloor$

\vspace{0.05cm}
\footnotesize
\begin{tabular}{@{}rrr@{}}
  \toprule
  \textbf{Tokens} & \textbf{Per sequence} & \textbf{Fit} \\
  \midrule
  8,192 & 2.7 GB & 195 \\
  32,768 & 10.7 GB & \textbf{48} \\
  131,072 & 42.9 GB & 12 \\
  \bottomrule
\end{tabular}
\end{column}
\begin{column}{0.50\textwidth}
\safeimg[width=\textwidth,height=4.3cm,keepaspectratio]{resident-fit-curve.pdf}
\end{column}
\end{columns}

\bottomline{Weights are paid once; every sequence adds private state, so resident fit falls as $1/S$.}
\end{frame}

\begin{frame}{Predict: what sets how many sequences a node holds?}
\note{
[3 min] Predict, then reveal.
% -- NARRATE
Give 60 seconds to write an answer and a reason. Most students will say
model size because the weights are the big number they know. Reveal: the
same model holds 195 sequences or 12, depending only on retained context.
% -- WARN
A smaller model raises capacity only through its per-token bytes or the
pool memory it frees. Parameter count is a poor predictor; the next slide
shows a 671B model that caches the least per token.
% -- FLEX
[CORE]
}

\centering
\vspace{0.4cm}
{\Large\bfseries Same node, same Llama 3 70B weights.}\\[0.4cm]
{\normalsize Which decides how many repair sequences fit at once:\\ the size of the model, or how many tokens each sequence retains?}\\[0.2cm]
{\small\textcolor{midgray}{Write your answer and one reason. 60 seconds.}}

\pause
\vspace{0.4cm}
\raggedright
\begin{mlsyscard}{crimson}
\small \textbf{Retained context.} The 526.0 GB pool holds \textbf{195} sequences at 8,192 tokens and \textbf{12} at 131,072, with the model unchanged. A smaller model helps only to the extent it lowers bytes per token or frees pool memory.
\end{mlsyscard}
\end{frame}

\begin{frame}{Attention design, fixed at pretraining, sets the slope}
\note{
[3 min]
% -- NARRATE
Multi-head attention stores every head; grouped-query attention lets several
query heads share one key-value head; multi-head latent attention stores one
compressed latent per token and rebuilds heads during attention. GPT-3's
geometry retains 4,718.6 KB per token, 67 times DeepSeek-V3's 70.3 KB.
Llama 3 70B, the reference model, retains 4.7 times DeepSeek-V3.
% -- ENGAGE
Ask: ``Can the serving engine drop half the key-value heads of a deployed
checkpoint to save memory?'' [No; head grouping is fixed at pretraining, and
dropping heads changes the computation.]
% -- WARN
Quantizing keys and values to 8 or 4 bits halves or quarters the element
bytes, but adds scale metadata and possible quality loss that must be
validated, which the approximation section returns to.
% -- FLEX
[OPTIONAL] Keep the 42.9 GB versus 9.2 GB contrast if short on time.
}

\begin{columns}[T]
\begin{column}{0.56\textwidth}
\safeimg[width=\textwidth,height=4.6cm,keepaspectratio]{kv-per-token-by-design.pdf}
\end{column}
\begin{column}{0.41\textwidth}
\small
\textbf{MHA} stores every head; \textbf{GQA} shares one KV head per query group; \textbf{MLA} stores one compressed latent per token.

\vspace{0.15cm}
GPT-3 geometry: \textbf{67$\times$} DeepSeek-V3 per token. Llama 3 70B: 4.7$\times$.

\vspace{0.15cm}
One 131,072-token session: \textbf{42.9 GB} (Llama 3 70B) vs \textbf{9.2 GB} (DeepSeek-V3).
\end{column}
\end{columns}

\bottomline{The checkpoint fixes bytes per token; the runtime's levers act only along that slope: tokens kept, sharing, precision, and eviction.}
\end{frame}

% =============================================================================
% SECTION 2: PAGING
% =============================================================================
\section{Paging}

\begin{frame}{Contiguous allocation strands memory in two ways}
\note{
[3 min]
% -- NARRATE
Left: a request admitted under a 32,768-token cap reserves 10.74 GB at once,
so the pool admits 48, whether or not the tokens arrive. If the population
finishes at a mean of 2,048 tokens, each reservation is 6.25\% full at
completion. Right: three free extents of 256 slots sit between live
allocations; 768 slots are free, yet a 512-slot contiguous request fails.
Growing a buffer instead means allocating a larger one elsewhere and copying,
holding both during the copy.
% -- ENGAGE
Ask: ``Is contiguity a hardware requirement?'' [No. The accelerator already
translates addresses; the attention kernel's fixed stride is what demands one
span. Software imposed it, so software can remove it.]
% -- FLEX
[CORE]
}

\centering
\safeimg[width=0.95\textwidth,height=2.15cm,keepaspectratio]{two-wastes.pdf}

\vspace{0.1cm}
\raggedright\footnotesize
\begin{tabular}{@{}L{0.20\textwidth}L{0.37\textwidth}L{0.36\textwidth}@{}}
  \toprule
  \textbf{Condition} & \textbf{Where capacity is lost} & \textbf{Cured by} \\
  \midrule
  Internal & reserved slots inside an allocation & not reserving for the cap \\
  External & free extents too small for a contiguous request & a uniform allocation unit \\
  \bottomrule
\end{tabular}

\vspace{0.1cm}
Early serving systems held token state in only \textbf{20.4--38.2\%} of allocated KV memory.

\bottomline{Full-cap reservation admits 48 r18-sized slots no matter how few tokens actually arrive.}
\end{frame}

\begin{frame}{A block table maps token positions to scattered blocks}
\note{
[3 min]
% -- LINK
Students know paged virtual memory from operating systems: the Atlas one-level
store introduced paging in 1962. The block table applies the same
indirection to token positions.
% -- NARRATE
Trace position 9 with four-token blocks: integer division gives logical block
2, the table maps it to physical block 9, and 9 mod 4 gives slot 1. Point out
the gray blocks: other requests' state sits between this sequence's blocks,
and it does not matter.
% -- WARN
The block table is software the attention kernel reads, not the
accelerator's hardware page table. And a mapped block's empty slots are
allocated capacity, not valid state.
% -- FLEX
[CORE]
}

\centering
\safeimg[width=0.90\textwidth,height=3.3cm,keepaspectratio]{block-table.pdf}

\vspace{0.05cm}
\raggedright\small
A \textbf{paged KV block table} maps a sequence's logical blocks to physical ones. All blocks are the same size, so any free block fits any request: no external fragmentation.

\bottomline{Indirection lets a sequence grow one block at a time without relocating what it already holds.}
\end{frame}

\begin{frame}{Paging bounds waste below one block and fits about 16$\times$ more}
\note{
[3 min]
% -- NARRATE
Tail waste: at most one block minus one slot per sequence. A 16-token block
is 5.24 MB across the node, a 64-token block 20.97 MB; even the larger tail is
0.2\% of a 32,768-token context. Then the payoff: final lengths with a mean of
2,048 under the 32,768 cap give about 780 final-length snapshots where full-cap
reservation gave 48.
% -- ENGAGE
Ask: ``Does r18, which actually grows to the cap, gain resident capacity from
paging?'' [No. The 16$\times$ is roughly cap over mean, 32,768 / 2,048. r18
still gains growth without relocation and block-level sharing.]
% -- WARN
The 780 is a planning estimate over completed requests, not a live admission
count: residents are mid-growth, and admission must reserve promised output.
% -- FLEX
[CORE]
}

\begin{columns}[T]
\begin{column}{0.48\textwidth}
\safeimg[width=\textwidth,height=2.8cm,keepaspectratio]{paging-payoff.pdf}\\[2pt]
{\scriptsize Final lengths with mean 2,048 tokens, capped at 32,768; snapshots at completion.}
\end{column}
\begin{column}{0.49\textwidth}
\footnotesize
\begin{tabular}{@{}lrr@{}}
  \toprule
  \textbf{Block size} & \textbf{16 tokens} & \textbf{64 tokens} \\
  \midrule
  Storage per block & 5.24 MB & 20.97 MB \\
  Worst tail (one block $-1$) & 4.92 MB & 20.64 MB \\
  Expected unused share & 0.4\% & 1.5\% \\
  \bottomrule
\end{tabular}

\vspace{0.15cm}
Tail waste $0 \le \Delta_{\text{tail}} \le \Sblk - 1$. Smaller blocks share more finely; larger blocks amortize bookkeeping and fit kernel tiles better.
\end{column}
\end{columns}

\bottomline{Paging's capacity gain measures how far requests fall short of the cap; at agent lengths, choose block size for kernels and sharing, not tail waste.}
\end{frame}

\begin{frame}{A forked candidate shares blocks instead of copying 10.7 GB}
\note{
[3 min]
% -- LINK
The breadth axis of test-time compute already priced parallel candidates in
memory: stored once, the prompt plus each candidate's own tokens. Paging is
the mechanism that delivers that.
% -- NARRATE
Four candidate repairs fork from one 8,192-token prefix. A contiguous
allocator writes four copies, 10.7 GB, the footprint of a whole r18 context.
A paged allocator clones four block tables and increments reference counts;
the work scales with block entries, not gigabytes.
% -- ENGAGE
Ask: ``When a candidate is discarded, when does its memory come back?''
[Each block returns when its reference count reaches zero and its pending
device operations finish.]
% -- FLEX
[CORE]
}

\centering
\safeimg[width=0.95\textwidth,height=3.05cm,keepaspectratio]{fork-share.pdf}

\vspace{0.1cm}
\raggedright\small
Physical memory is the count of \textbf{unique} blocks resident, not the sum of the candidates' logical contexts.

\bottomline{Under paging, a fork costs block-table entries; only the tokens each candidate writes cost new memory.}
\end{frame}

\begin{frame}{Copy-on-write copies only a shared, partly filled tail}
\note{
[3 min]
% -- NARRATE
Both candidates share table [3, 7, 9]. Block 9 holds positions 8 and 9.
Before B writes position 10, the engine allocates block 14, copies the two
initialized entries, points B's last entry at 14, and B writes into slot 2.
A is untouched. Copy-on-write comes from TENEX in the early 1970s.
% -- ENGAGE
Ask: ``The prefix is 12 tokens instead of 10. Is any copy needed?'' [No.
Position 12 starts a new logical block, so the candidate allocates a fresh
private block.]
% -- WARN
Students think every fork triggers copying. Only an append into a shared,
partly filled block does.
% -- FLEX
[CORE]
}

\centering
\safeimg[width=0.92\textwidth,height=2.8cm,keepaspectratio]{cow-tail.pdf}

\vspace{0.05cm}
\raggedright\footnotesize\renewcommand{\arraystretch}{0.9}
\begin{tabular}{@{}L{0.33\textwidth}L{0.60\textwidth}@{}}
  \toprule
  \textbf{Append target} & \textbf{Engine action} \\
  \midrule
  shared, partly filled block & copy initialized entries to a private block, then write \\
  next position starts a new block & allocate a fresh private block; no copy \\
  private block (one owner) & write into the remaining slot in place \\
  \bottomrule
\end{tabular}

\bottomline{Two state variables decide every append: how full the target block is and how many owners it has.}
\end{frame}

\begin{frame}{Your turn: what does a longer mean change?}
\note{
[3 min] Exercise, from the chapter's allocation checkpoint.
% -- NARRATE
Give 90 seconds. Answers: 4,096 / 32,768 is 12.5\% populated at completion.
Static reservation still admits 48, because a reservation is sized by the
cap, not by the tokens that arrive. Paging's snapshot gain is roughly cap
over mean, 32,768 / 4,096 = 8, so about half the earlier advantage.
% -- WARN
The 8$\times$ is the chapter's cap-over-mean rule applied to new inputs, a
planning estimate, not a measured admission count.
% -- FLEX
[OPTIONAL] Skip if the paging payoff slide ran long.
}

\centering
\vspace{0.2cm}
{\large The mean final length doubles from 2,048 to \textbf{4,096} tokens.\\ The cap stays at 32,768 tokens.}\\[0.35cm]
\raggedright\small
\begin{enumerate}\setlength{\itemsep}{3pt}
  \item What fraction of each full-cap reservation holds state at completion?
  \item Does static reservation now admit more than 48 sequences?
  \item Roughly how large is paging's capacity advantage now?
\end{enumerate}

\pause
\begin{mlsyscard}{crimson}
\small (1) $4{,}096/32{,}768 = $ \textbf{12.5\%}. (2) \textbf{No}: reservations are sized by the cap, so still 48. (3) About cap over mean, $32{,}768/4{,}096 \approx$ \textbf{8$\times$}, half the earlier gain.
\end{mlsyscard}
\end{frame}

% =============================================================================
% SECTION 3: PREFIX CACHE
% =============================================================================
\section{Prefix Cache}

\begin{frame}{A prefix cache finds reusable state across requests}
\note{
[2 min]
% -- LINK
Context engineering priced a reused prefix as a discounted cached-input
line on the bill. Today is what produces that discount inside the engine.
% -- NARRATE
The turn that processes the test result starts with the same system prompt,
tools, and history as the turn before. If the engine kept that state, it
attaches it and prefills only the new tokens. Contrast with forking: the
engine forked the candidates itself and knows they are related; a prefix
cache must discover reusable state, by content, among requests that arrive
independently, and can offer only what eviction left behind.
% -- FLEX
[CORE]
}

\begin{columns}[T]
\begin{column}{0.52\textwidth}
\begin{mlsyscard}{bookblue}
\small\textbf{Prefix cache:} the serving system's index of retained KV state, keyed by exact token prefix and computation configuration, so a later request sharing that prefix attaches existing state instead of recomputing it.
\end{mlsyscard}
\end{column}
\begin{column}{0.44\textwidth}
\footnotesize
\begin{tabular}{@{}L{0.30\textwidth}L{0.60\textwidth}@{}}
  \toprule
  & \textbf{How state is found} \\
  \midrule
  Block sharing & engine forked the requests; relation known \\
  Prefix cache & matched by content across independent requests \\
  \bottomrule
\end{tabular}
\end{column}
\end{columns}

\vspace{0.25cm}
\small
Three parts make it a system rather than a discount: an \textbf{index} over retained prefixes, an \textbf{ownership contract} on the blocks it hands out, and an \textbf{eviction policy} under memory pressure.

\bottomline{On the bill a hit is a discount; in the engine it is a lookup, a reference count, and a retention decision.}
\end{frame}

\begin{frame}{Reuse requires compatible, available, permitted state}
\note{
[3 min]
% -- NARRATE
Compatible: an exact, contiguous token prefix from position zero, because
attention state is derived from the token prefix, and the same weights,
adapters, positional encodings, masks, and any multimodal inputs. Available:
compatible tensors may already be evicted, or sit in host memory and need a
page-in. Permitted: resident tensors may belong to another tenant. Point to
the dashed box: no gate asks whether the evidence is still true.
% -- ENGAGE
Ask: ``Two prompts say the same thing in different words. Reusable?'' [No.
Only identical tokens from the first position match.]
% -- WARN
The reusable length is often shorter than the longest match in the tree, and
new tokens join the index only after their blocks are committed.
% -- FLEX
[CORE]
}

\centering
\safeimg[width=0.95\textwidth,height=2.65cm,keepaspectratio]{reuse-gates.pdf}

\vspace{0.1cm}
\raggedright\small
\textbf{Reusable prefix length} = the longest contiguous prefix that passes all three gates. \textbf{Compatible} means identical tokens from position 0 \emph{and} identical weights, adapters, positional encodings, masks, and multimodal inputs.

\bottomline{A prefix is reusable only when it is compatible, still available, and permitted, and none of the three says whether its evidence is current.}
\end{frame}

\begin{frame}{A radix tree returns the longest retained matching prefix}
\note{
[3 min]
% -- LINK
Context staging ordered the prompt by change rate. Here is why that order
pays: every position's KV depends on the full history before it.
% -- NARRATE
Each node is a run of tokens and the blocks that hold them. Walk from the
root until the request diverges: that one walk returns both the matched
length and the blocks. The engine increments reference counts on those
blocks before eviction can reclaim them. Lookup costs microseconds on the
host. A one-token-per-edge trie would spend a node on each of thousands of
prefix tokens; the radix tree compresses chains into single edges.
% -- ENGAGE
Ask: ``A timestamp is serialized after token 300. What does the next turn
reuse?'' [At most the first 300 tokens.]
% -- FLEX
[CORE]
}

\begin{columns}[T]
\begin{column}{0.60\textwidth}
\safeimg[width=\textwidth,height=4.6cm,keepaspectratio]{radix-lookup.pdf}
\end{column}
\begin{column}{0.37\textwidth}
\small
\textbf{Lookup:} walk edges from the root to the divergence point; return the match length and its blocks.

\vspace{0.1cm}
\textbf{Layout:} volatile content (counters, tool outputs, notes) last keeps the invariant prefix on one path.

\vspace{0.1cm}
\textbf{Block boundary:} a 10-token match with 4-token blocks reuses 10 with partial tails, 8 with full blocks only.
\end{column}
\end{columns}

\bottomline{A node's KV state is valid only beneath its exact ancestors, so one changed early token forfeits everything after it.}
\end{frame}

\begin{frame}{Leaf-first eviction keeps residency prefix-closed}
\note{
[3 min]
% -- NARRATE
Under pressure, only dormant leaves are eligible: here the rejected branch
and the final turn. Other task is dormant but has a child, so it waits;
once final turn is evicted, other task becomes an eligible leaf. Within the
eligible set, engines often pick least recently used. Prefix-closed means
every resident node's ancestors are resident: a suffix without its prefix
is useless state in a causal model.
% -- ENGAGE
Ask: ``A system prompt shared by hundreds of workers goes quiet for a minute.
Should LRU evict it?'' [Recency says yes; its reuse value says no. Better
schemes weigh reuse probability and rebuild cost against footprint.]
% -- WARN
Belady's optimal policy needs the future. An agent harness sometimes holds
part of it, such as a tool launch that will idle a prefix for minutes, so a
caller's retention hint can beat recency alone.
% -- FLEX
[CORE]
}

\begin{columns}[T]
\begin{column}{0.60\textwidth}
\safeimg[width=\textwidth,height=4.6cm,keepaspectratio]{leaf-first-eviction.pdf}
\end{column}
\begin{column}{0.37\textwidth}
\small
\textbf{Eligible:} dormant leaves only; least recently used among them.

\vspace{0.1cm}
\textbf{Cascade:} a parent becomes eligible once its last child is evicted.

\vspace{0.1cm}
\textbf{Why:} evicting an ancestor would orphan every descendant, which then needs a reload or a new prefill before it is usable.
\end{column}
\end{columns}

\bottomline{An ancestor never leaves before its descendants, because a suffix without its prefix cannot be used.}
\end{frame}

\begin{frame}{Discussion: does a cache hit prove the evidence is current?}
\note{
[4 min] Discussion in pairs.
% -- LINK
Context invalidation, from context engineering, is how the harness
withdraws a stale observation while keeping it as history.
% -- NARRATE
Two minutes in pairs, then reveal. The r17 tensors are numerically exact for
the r17 prompt; nothing about them is wrong. What is wrong is the claim they
support. Freshness is enforced by the prompt layout: the harness writes the
r18 observation, the match ends at the first changed token, and the turn
pays a partial prefill.
% -- WARN
Some students will propose ``invalidate the cache when the repository
changes.'' The cache cannot know what the tokens mean; only the harness can.
% -- FLEX
[CORE]
}

\centering
\vspace{0.1cm}
{\large The cached prefix contains the \textbf{r17} test failure. The repository is now at \textbf{r18}.}\\[0.2cm]
{\normalsize The tokens match, the blocks are resident, and the tenant is permitted.}\\[0.25cm]
\raggedright\small
\begin{enumerate}\setlength{\itemsep}{2pt}
  \item Is the cached KV state valid? Is the evidence in it current?
  \item Who must enforce freshness, and what does enforcing it cost?
\end{enumerate}

\pause
\begin{mlsyscard}{crimson}
\small The state is exact for the r17 prompt, but \textbf{reusable is not applicable}. The \textbf{harness} rewrites the prompt with r18 evidence; the match stops at the first changed token, forfeiting part of a 169 ms saving. That is the correct price of grounding the next decision in its source.
\end{mlsyscard}
\end{frame}

% =============================================================================
% SECTION 4: PREFILL
% =============================================================================
\section{Prefill}

\begin{frame}{A hit replaces 169 ms of prefill with a 5 ms pass}
\note{
[3 min]
% -- LINK
Prefill, decode, and time to first token come from the call-latency
decomposition earlier in the book.
% -- NARRATE
A deliberately small 4,096-token invariant prefix: 1.34 GB of KV. Rebuilding
it takes about 578 trillion operations for projections plus 22 trillion for
attention. At 45\% of the node's 16-bit peak, about 3,560 trillion per
second, that is 169 ms. A hit that appends 32 tokens is one pass bounded by
reading the weights once: about 5 ms. The longer shared zone of r18 would
save more.
% -- ENGAGE
Ask: ``Why is the 32-token pass not 32/4,096 of 169 ms?'' [Too little
arithmetic to fill the matrix units; the weight read sets a floor.]
% -- FLEX
[CORE]
}

\begin{columns}[T]
\begin{column}{0.52\textwidth}
\safeimg[width=\textwidth,height=2.8cm,keepaspectratio]{prefix-hit-vs-miss.pdf}

\vspace{0.1cm}
\footnotesize
\textbf{Step-time model} for a pass over $n$ positions:\\
$T(n) = \max\big(T_{\text{read}},\; 2Pn/\Rattain\big)$, weight-read floor about 5 ms.
\end{column}
\begin{column}{0.45\textwidth}
\small
\textbf{Miss:} rebuild 4,096 tokens (1.34 GB of KV): 578 T operations for projections + 22 T for attention at about 3,560 T operations/s $\approx$ \textbf{169 ms}.

\vspace{0.15cm}
\textbf{Hit:} attach the prefix, run the 32-token suffix: one weight-bound pass, about \textbf{5 ms}.
\end{column}
\end{columns}

\bottomline{A trajectory that keeps its invariant prefix stable collects this saving on every turn, in latency and in total prefill work.}
\end{frame}

\begin{frame}{Measure reuse by tokens actually served from cached state}
\note{
[3 min]
% -- NARRATE
The numerator counts only tokens attached after availability, permission,
and block rounding: the 10-token match under full-block caching counts 8.
For a hosted API, the cached-token field in the usage record is that
numerator for one request. Track it across turns. A plunge to zero after a
long tool wait is consistent with idle eviction, the remote form of drop and
recompute, but memory pressure and routing produce the same signal. A count
stuck below expectation points to an early serialization difference.
% -- ENGAGE
Ask: ``The cached count is stuck at 300 of 6,000 tokens every turn. First
suspect?'' [Something volatile serialized near token 300, such as a
timestamp.]
% -- WARN
Hit rate alone can mislead: a hit removes prefix compute, and then suffix
prefill, decode, or queueing can dominate latency.
% -- FLEX
[CORE]
}

\centering
{\large $\rho_{\text{prefix}} = \dfrac{\sum_i S_{\text{match}}^{(i)}}{\sum_i S_{\text{prompt}}^{(i)}}$}
\quad{\footnotesize after availability, permission, and block rounding}

\vspace{0.15cm}
\raggedright\footnotesize
\begin{tabular}{@{}L{0.37\textwidth}L{0.56\textwidth}@{}}
  \toprule
  \textbf{Caller observes} & \textbf{Consistent with} \\
  \midrule
  cached count drops to zero after a long tool wait & idle eviction (or pressure, routing, changed settings) \\
  cached count stable but below expectation & early serialization divergence: timestamps, whitespace, reordered JSON keys \\
  \bottomrule
\end{tabular}


\bottomline{The serialized context, not the cache, is the model's input: a miss changes latency, never correctness.}
\end{frame}

\begin{frame}{A monolithic prefill stalls every ongoing decode}
\note{
[3 min]
% -- LINK
Continuous batching re-forms the batch at every step. That per-iteration
boundary is where a slice of one request's prefill can join the decodes.
% -- NARRATE
Thirty-two requests decode at the weight-read floor, about 5 ms per step. A
4,096-token prompt admitted whole adds about 162 ms of weight arithmetic, so
the decodes wait about 168 ms between tokens. Split into 512-token chunks,
each iteration carries 544 positions, about 22 ms, and the decodes advance
every iteration. The new prompt finishes at about 173 ms instead of 162.
% -- ENGAGE
Ask: ``Who pays for chunking, and how much?'' [The new request, about 11 ms
of time to first token; the 32 streams gain an 8$\times$ smaller worst gap.]
% -- FLEX
[CORE]
}

\centering
\safeimg[width=0.95\textwidth,height=3.9cm,keepaspectratio]{chunked-timeline.pdf}

\vspace{0.05cm}
\raggedright\footnotesize
\textbf{Chunked prefill} splits an uncached prompt into spans, each run beside the ongoing decodes. Weight-only estimate: attention and overhead lengthen every bar.

\bottomline{Chunking cuts the worst gap between tokens about 8$\times$ for an 11 ms delay to the new request's first token.}
\end{frame}

\begin{frame}{Pick the largest chunk that meets the token-gap target}
\note{
[3 min]
% -- NARRATE
Each decode contributes one position, each prefill chunk its length.
Decode-first scheduling reserves the decode slots, then fills the rest with
prefill. Operators set the budget to the largest chunk whose iteration
comfortably meets the inter-token target: on the reference node, a target
comfortably above 22 ms admits 512-token chunks beside 32 decodes. Then
measure: attention traffic and overhead shrink that chunk.
% -- ENGAGE
Ask: ``Why judge by the p99 inter-token gap, not the mean?'' [The mean hides
the occasional large prefill stall that users feel.]
% -- WARN
Disaggregation puts prefill and decode on separate pools instead, but every
prompt's KV must then cross to a decode worker; when that link saturates it
can cost more than isolation saves.
% -- FLEX
[OPTIONAL]
}

\centering
{\large $\displaystyle\sum_i C_i + |\mathcal{D}| \le \Nbudget$}\quad{\footnotesize prefill chunk positions + one per decode}

\vspace{0.15cm}
\raggedright\footnotesize
\begin{tabular}{@{}L{0.20\textwidth}L{0.36\textwidth}L{0.36\textwidth}@{}}
  \toprule
  & \textbf{Smaller chunks} & \textbf{Larger chunks} \\
  \midrule
  Ongoing streams & short, frequent gaps between tokens & longer stalls \\
  New request & later first token; more launches & earlier first token \\
  Accelerator & narrower matrices, more overhead & wider matrices, better utilization \\
  \bottomrule
\end{tabular}

\vspace{0.1cm}
Calibrate by replaying arrival traces and reading time to first token, \textbf{p99} inter-token gap, throughput, and memory occupancy.

\bottomline{The budget is the largest chunk whose measured iteration still meets the inter-token target; the weight-only estimate is only the starting point.}
\end{frame}

% =============================================================================
% SECTION 5: TOOL WAITS
% =============================================================================
\section{Tool Waits}

\begin{frame}{At a tool wait, the engine keeps, moves, or drops the state}
\note{
[3 min]
% -- LINK
Invariant closure puts mechanical bounds below the model. The KV pool is
such a bound.
% -- NARRATE
Read the table by rows: keep releases nothing and needs nothing on return;
move releases blocks after a verified copy and reloads them; drop releases
after safe release and rebuilds from the recorded tokens. Then the right
card: the model cannot see the pool, cannot enforce a release, and the next
turn's correctness does not depend on where the state was kept. So the
engine decides; the agent runtime may only hint.
% -- WARN
Students assume the agent should manage its own memory. The model neither
makes the decision nor supplies the hints.
% -- FLEX
[CORE]
}

\footnotesize
\begin{tabular}{@{}L{0.09\textwidth}L{0.26\textwidth}L{0.26\textwidth}L{0.26\textwidth}@{}}
  \toprule
  \textbf{Policy} & \textbf{Memory released} & \textbf{Work on return} & \textbf{Constraint} \\
  \midrule
  \textcolor{bookgreen}{\textbf{Keep}} & none & nothing & competes for memory \\
  \textcolor{bookviolet}{\textbf{Move}} & after a verified copy & allocate, reload & host space, links \\
  \textcolor{bookorange}{\textbf{Drop}} & after safe release & rebuild from inputs & prefill capacity \\
  \bottomrule
\end{tabular}

\vspace{0.2cm}
\begin{mlsyscard}{bookviolet}
\footnotesize\textbf{Who decides:} the \textbf{serving engine}. The model cannot see the pool or enforce a release, and the next turn is correct wherever the state was kept. The \textbf{agent runtime} may hint the expected wait.
\end{mlsyscard}

\bottomline{The pause policy is a serving decision below the model, informed by hints from the agent runtime.}
\end{frame}

\begin{frame}{Only unshared, idle blocks can actually be freed}
\note{
[3 min]
% -- NARRATE
Start every pause calculation from the uniquely reclaimable footprint: shared
prefix blocks stay while another request references them, and blocks in an
in-flight copy or kernel stay until it completes. Moving keeps the source
blocks until the outbound copy verifies, so during the transfer the state
occupies device and host memory at once. Dropping keeps the information as
long as the exact token record, adapters, positions, masks, and multimodal
inputs are retained; a summary cannot rebuild it.
% -- LINK
Durable execution owns the trajectory log; here we only need that the log,
not the KV cache, is the record.
% -- FLEX
[CORE]
}

\begin{columns}[T]
\begin{column}{0.48\textwidth}
\small
\textbf{Reclaimable footprint} excludes
\begin{itemize}\setlength{\itemsep}{1pt}
  \item prefix blocks another request still references
  \item blocks under an in-flight copy or kernel
\end{itemize}
\textbf{Move} holds source blocks until the copy verifies: for that window the state occupies both device and host memory. On resume: allocate, transfer, update tables.
\end{column}
\begin{column}{0.48\textwidth}
\begin{mlsyscard}{bookorange}
\small\textbf{Drop} forfeits computation, not information, if the exact inputs are kept: token sequence, adapters, positional encodings, masks, multimodal inputs.\\[3pt]
A conversation summary cannot rebuild the state. Rebuild reuses any compatible resident prefix and prefills only the rest.
\end{mlsyscard}
\end{column}
\end{columns}

\bottomline{KV state is derived and recomputable; the runtime's event log, not the cache, is the record of the task.}
\end{frame}

\begin{frame}{Reloading paused state beats recomputing it by about 62$\times$}
\note{
[3 min]
% -- NARRATE
Each of 8 devices holds an eighth of the state and has its own 64 GB/s
link: 512 GB/s together. Moving 10.74 GB takes 21 ms each way. Recompute
is weight work, $2PS$ operations at 3,560 trillion per second: 1.3 s. Both
scale linearly with $S$, so the ratio does not depend on context length.
Even at a quarter of the link rate, reload is 15$\times$ faster.
% -- ENGAGE
Ask: ``When would recompute win?'' [When host memory or the links are
contested and accelerator compute is the idle resource.]
% -- WARN
The incoming test observation must be prefilled under every policy, so it
is left out of the comparison.
% -- FLEX
[CORE]
}

\centering
\safeimg[width=0.95\textwidth,height=2.8cm,keepaspectratio]{reload-vs-recompute.pdf}

\vspace{0.1cm}
{\large $\dfrac{T_{\text{reload}}}{T_{\text{recompute}}} = \dfrac{S\,\mtok/\BWlink}{2PS/\Rattain} = \dfrac{\mtok\,\Rattain}{2P\,\BWlink}$}

\vspace{0.05cm}
\raggedright\footnotesize
8 host links $\times$ 64 GB/s $=$ 512 GB/s; recompute of $2PS$ weight operations at about 3,560 T operations/s (attention would add more).

\bottomline{On this node, moving state back is cheap and rebuilding it is not, at every context length.}
\end{frame}

\begin{frame}{Your turn: price the r18 test wait}
\note{
[4 min] Exercise.
% -- NARRATE
Give two minutes. Keep: nothing released, 10.74 GB for 45 s is 483 GB-seconds
of memory-time, no resume delay. Move: blocks free after the 21 ms outbound
copy, for 44.98 s; the 21 ms reload lands on the resume. Drop: memory free
for the whole wait; 1.3 s of prefill on the resume. Move with prefetch timed
to the test's end: free for 44.96 s and no delay, if the return time is known.
% -- WARN
Students count the 42 ms round trip as resume delay. Only the inbound reload
sits on the critical path.
% -- FLEX
[CORE]
}

\small
r18 holds \textbf{10.74 GB}, all reclaimable, through a \textbf{45 s} test. Transfer: 21 ms each way. Rebuild: 1.3 s.\\[2pt]
For each policy: \textbf{how long is the memory free}, and \textbf{what delay lands on the resume}?

\pause
\vspace{0.1cm}
\footnotesize
\begin{tabular}{@{}L{0.24\textwidth}L{0.33\textwidth}L{0.33\textwidth}@{}}
  \toprule
  \textbf{Policy} & \textbf{Accelerator memory} & \textbf{Delay on resume} \\
  \midrule
  Keep & held: 483 GB$\cdot$s of memory-time & none \\
  Move & free for 45 s $-$ 21 ms $=$ 44.98 s & 21 ms reload \\
  Drop & free for the whole wait & 1.3 s rebuild \\
  Move + prefetch & free for 44.96 s & none, if the return time is known \\
  \bottomrule
\end{tabular}

\bottomline{For a long wait under memory pressure, moving reclaims almost the whole wait for a 21 ms resume.}
\end{frame}

\begin{frame}{The wait, relative to the round trip, picks the regime}
\note{
[3 min]
% -- NARRATE
The diagonal is the round trip, twice the state over the link rate: 42 ms at
32,768 tokens, growing with context. A wait shorter than that ends before
the state can go out and back, so keep it if memory permits. Above it,
moving is possible, and it pays when other requests want the memory. Far to
the right, parked state piles up in host memory; the dashed line depends on
host capacity and on how fast paused trajectories arrive, not a fixed time.
r18 sits deep in the move regime.
% -- ENGAGE
Ask: ``The pool has plenty of free blocks. Should r18 still move?'' [No:
42 ms of transfer and 21 ms on resume buy nothing.]
% -- FLEX
[CORE]
}

\begin{columns}[T]
\begin{column}{0.58\textwidth}
\safeimg[width=\textwidth,height=4.5cm,keepaspectratio]{pause-regimes.pdf}
\end{column}
\begin{column}{0.39\textwidth}
\small
\textbf{Round trip} $= 2\Mkv/\BWlink$: 42 ms for r18.

\vspace{0.1cm}
\textbf{Wait $<$ round trip:} keep if memory permits.

\vspace{0.1cm}
\textbf{Wait $>$ round trip:} move pays when other requests wait for memory.

\vspace{0.1cm}
\textbf{Minutes to hours} (human approval): host memory fills; park on storage or drop.
\end{column}
\end{columns}

\bottomline{A tool call is an opportunity for a scheduling decision, not an automatic eviction trigger.}
\end{frame}

\begin{frame}{Correlated returns turn a 21 ms reload into 0.67 s}
\note{
[3 min]
% -- NARRATE
Thirty-two paused repairs finish their tests together: 344 GB to reload, at
least 0.67 s even at the full 512 GB/s. The last one waits 32$\times$ its
isolated 21 ms, longer if shared bus root complexes congest. Then the three
ways return-time variance bites: early return during the outbound copy,
premature prefetch that re-occupies memory while the model idles, and a
cancelled trajectory whose reload was wasted.
% -- ENGAGE
Ask: ``How can a reload be hidden?'' [Stream layer by layer, reloading layer
$\ell+1$ while attention runs on layer $\ell$, or run unrelated requests
during the reload; both need reserved headroom.]
% -- FLEX
[OPTIONAL]
}

\begin{columns}[T]
\begin{column}{0.50\textwidth}
\safeimg[width=\textwidth,height=4.3cm,keepaspectratio]{return-burst.pdf}
\end{column}
\begin{column}{0.47\textwidth}
\small
\textbf{32 returns at once:} $32 \times 10.74 \approx 344$ GB, $\ge$ \textbf{0.67 s} at 512 GB/s.

\vspace{0.1cm}
\textbf{Return time varies:}
\begin{itemize}\setlength{\itemsep}{1pt}
  \item early return during the outbound copy
  \item early prefetch re-occupies memory
  \item a cancelled trajectory wastes its reload
\end{itemize}
\end{column}
\end{columns}

\bottomline{Plan pause policy against the distribution of return times and burst reloads, not the mean wait.}
\end{frame}

% =============================================================================
% SECTION 6: APPROXIMATION
% =============================================================================
\section{Approximation}

\begin{frame}{Lossy approximation changes the computation itself}
\note{
[3 min]
% -- NARRATE
Keep, move, and drop all preserve full attention exactly, up to
floating-point differences. A second family saves memory by approximating:
evict tokens, attend only to a recent window, or quantize keys and values.
Attention sinks show the danger: softmax must put its mass somewhere, and
layers deposit it on the first positions. A sliding window that drops them
makes quality collapse past the window; keeping a few sink tokens restores
stability.
% -- WARN
Students read memory saved as success. A lossy policy's benefit is
measured, but its cost is a changed computation that only task-level
evaluation reveals.
% -- FLEX
[CORE]
}

\footnotesize
\begin{tabular}{@{}L{0.17\textwidth}L{0.36\textwidth}L{0.38\textwidth}@{}}
  \toprule
  & \textbf{Lossless: keep, move, drop} & \textbf{Lossy: evict, window, quantize} \\
  \midrule
  Computation & full attention, same inputs & changed: fewer or coarser KV \\
  Pays with & link bandwidth, tiered memory, prefill & task accuracy, which must be measured \\
  Judged by & memory freed, resume latency & the task's hard dependencies \\
  \bottomrule
\end{tabular}

\vspace{0.2cm}
\begin{mlsyscard}{bookred}
\footnotesize\textbf{Attention sinks.} Softmax weights sum to one, so layers park leftover attention on the first positions. A sliding window that evicts them degrades sharply once context exceeds the window; keeping a few initial sink tokens restores stability. Score-based policies instead keep tokens with high accumulated or recent attention.
\end{mlsyscard}

\bottomline{Memory saved is not evidence: a lossy policy ships only after validation against the task.}
\end{frame}

\begin{frame}{Agent trajectories break ``low attention means irrelevant''}
\note{
[3 min]
% -- NARRATE
Score-based eviction assumes that what was ignored will stay ignored. Three
agent patterns violate it. A security rule set in the system prompt is
ignored for dozens of turns, then decides the final edit. A commit hash
appears once and must later be quoted exactly in a shell call. Brackets and
colons in a tool schema get little attention per step, yet every tool call's
grammar depends on them.
% -- ENGAGE
Ask: ``Which failure would perplexity catch?'' [None of them: each looks
like fluent text until a parser, a command, or a policy check fails.]
% -- FLEX
[CORE]
}

\begin{columns}[T]
\begin{column}{0.32\textwidth}
\begin{mlsyscard}{bookblue}
\footnotesize\textbf{Dormant constraints}\\[2pt]
Environment rules and security policies set early get little attention for dozens of turns, then govern the final edit. Pruned: silent task failure.
\end{mlsyscard}
\end{column}
\begin{column}{0.32\textwidth}
\begin{mlsyscard}{bookviolet}
\footnotesize\textbf{Exact identifiers}\\[2pt]
Commit hashes, process IDs, file paths, and nonces appear once, then go unreferenced. Pruned: the agent cannot quote them in a later call.
\end{mlsyscard}
\end{column}
\begin{column}{0.32\textwidth}
\begin{mlsyscard}{bookorange}
\footnotesize\textbf{Tool-schema syntax}\\[2pt]
Brackets, colons, and delimiters get weak attention per step but define every call's grammar. Pruned: malformed calls that crash parsers.
\end{mlsyscard}
\end{column}
\end{columns}

\bottomline{Judge approximation against hard semantic dependencies, not aggregate perplexity: past attention does not predict an agent's future needs.}
\end{frame}

\begin{frame}{Discussion: quantize the cache, or cap the context?}
\note{
[4 min] Discussion in pairs.
% -- LINK
Choosing the working set is context engineering's job; here we only weigh it
against approximation.
% -- NARRATE
Two minutes in pairs. Expected: both shrink the footprint, but they move
different things. Capping retained context keeps the computation exact over
what remains and makes the harness choose what the model sees, so the
evidence is task success with the trimmed working set. Quantizing changes
the computation for every token, so the evidence is accuracy on dormant
constraints, exact identifiers, and schema syntax, not perplexity. Press
each pair to name the measurement its proposal needs.
% -- FLEX
[CORE] Exercise 1(c) in the chapter quantifies this trade.
}

\centering
\vspace{0.2cm}
{\large An operator must hold more repair sessions resident than the pool allows.}\\[0.3cm]
{\normalsize Proposal A: store keys and values at \textbf{one byte} per element.\\ Proposal B: the harness \textbf{caps retained context} per session.}\\[0.35cm]
\raggedright\small
\begin{enumerate}\setlength{\itemsep}{3pt}
  \item What does each proposal change: the computation, or what the model sees?
  \item What evidence must each produce before it ships? Which would you try first?
\end{enumerate}
\centering
{\small\textcolor{midgray}{Two minutes with a neighbor.}}
\end{frame}

% =============================================================================
% SECTION 7: CAPACITY
% =============================================================================
\section{Capacity}

\begin{frame}{Sharing a block-aligned prefix lifts resident fit 3.4$\times$}
\note{
[4 min]
% -- NARRATE
Now plan honestly: reserve each request's 2,048-token output allowance,
withhold 10\% headroom, round to 16-token blocks of 5.243 MB. The pool gives
90,292 usable blocks; an unshared request needs 2,176, so 41 fit, down from
48. A fleet of repair workers on one repository snapshot shares its first
24,576 tokens: store those 1,536 blocks once and each request needs only 640
private blocks, so 138 fit. The shared fraction is 70.6\%, and $1/(1-f)$
predicts 3.4$\times$.
% -- ENGAGE
Ask: ``If the shared prefix shrank to 12,288 tokens, would sharing still
more than double the fit?'' [Shared fraction about 768/2,176, 35\%, gain about
1.5$\times$: no.]
% -- WARN
Sharing does not reduce decode work per token or change any request's
context; it multiplies concurrent capacity only.
% -- FLEX
[CORE]
}

\centering\small
$\lceil \Spre/\Sblk \rceil + B\,\lceil (S + \Kmax - \Spre)/\Sblk \rceil \;\le\; \lfloor (1-\fhead)\,\Mpool / (\Sblk\,\mtok) \rfloor$

\raggedright
\begin{columns}[T]
\begin{column}{0.50\textwidth}
\footnotesize
32,768 in + 2,048 out; 10\% headroom; 16-token blocks; a 24,576-token prefix shared:\\[2pt]
\begin{tabular}{@{}lr@{}}
  \toprule
  Usable blocks & 90,292 \\
  Unshared: 2,176 per request & \textbf{41} fit \\
  Shared prefix: 1,536 once & \\
  \quad + 640 private per request & \textbf{138} fit \\
  \bottomrule
\end{tabular}
\end{column}
\begin{column}{0.47\textwidth}
\safeimg[width=\textwidth,height=3.9cm,keepaspectratio]{sharing-gain.pdf}
\end{column}
\end{columns}

\bottomline{In a large pool the sharing gain approaches $1/(1-f)$, so prompt layout and the scope of sharing are the levers.}
\end{frame}

\begin{frame}{Little's law turns a slot count into a turn rate}
\note{
[4 min] Exercise at the end.
% -- NARRATE
Little's law holds for any stable system whatever the arrival pattern or
order. Read it backward: a fixed number of slots sustains a rate equal to
slots over residence. Residence counts every phase the slot is allocated,
the tool wait included. Then pose the question and give 90 seconds. Answer:
41 over 50 s is 0.82 turns per second; 41 over 5 s is 8.2; the ratio is
1 + 45/5 = 10.
% -- WARN
The law fixes the mean occupancy only. A burst of arrivals, or residences
that end together, holds far more slots than the mean while it lasts.
% -- FLEX
[CORE]
}

\begin{columns}[T]
\begin{column}{0.47\textwidth}
\centering
{\Large $\Nslot = \lam\,\Thold$}\\[0.15cm]
\raggedright\footnotesize
$\Nslot$: mean slots held \quad $\lam$: turns per second\\
$\Thold$: mean time a slot stays allocated, \textbf{waits included}

\vspace{0.15cm}
\small Rearranged: $\lam = \Nslot / \Thold$.\\ The law fixes the mean, not the peak.
\end{column}
\begin{column}{0.49\textwidth}
\small
\textbf{Your turn.} The node holds \textbf{41} slots. Each repair turn generates for \textbf{5 s}, then waits \textbf{45 s} for its test. What turn rate can memory sustain if state is kept through the wait? If the slot is released?

\pause
\vspace{0.1cm}
\begin{mlsyscard}{crimson}
\small Kept: $41/50\text{ s} = $ \textbf{0.82 turns/s}.\\ Released: $41/5\text{ s} = $ \textbf{8.2 turns/s}.\\ Ratio $1 + 45/5 = $ \textbf{10$\times$}.
\end{mlsyscard}
\end{column}
\end{columns}

\bottomline{A slot count is a concurrency bound; dividing it by residence time gives the sustainable rate.}
\end{frame}

\begin{frame}{A kept slot is paid for across the whole turn cycle}
\note{
[3 min]
% -- LINK
Cost per accepted task was defined in the introduction; this same factor
scales the memory term of every attempt's cost.
% -- NARRATE
The residency tax is $1 + \Ttool/\Tmodel$: what matters is the ratio of
waiting to generating, not the wait alone. Repair pays 10$\times$, support
11$\times$. The released rate is an ideal ceiling from memory alone: prefill
compute, restoration bandwidth, and scheduling often saturate first, and
shared prefix blocks stay pinned regardless.
% -- ENGAGE
Ask: ``The test doubles to 90 s; generation stays 5 s. What happens?'' [The
kept rate falls to 41/95 s, about 0.43 turns/s; the case for moving gets
stronger.]
% -- FLEX
[CORE]
}

\centering
\safeimg[width=0.95\textwidth,height=2.7cm,keepaspectratio]{residency-tax.pdf}

\vspace{0.1cm}
{\large Residency tax $= 1 + \dfrac{\Ttool}{\Tmodel}$}

\vspace{0.05cm}
\raggedright\footnotesize
Released-slot rates are memory-only ceilings: prefill, restoration bandwidth, and queues often bind first, and bursts push peaks above the mean.

\bottomline{A slot held through an idle wait divides the memory-bound turn rate by $1 + \Ttool/\Tmodel$.}
\end{frame}

\begin{frame}{Who waits on the resume decides the pause policy}
\note{
[3 min]
% -- LINK
The four lighthouse workloads from the introduction: here the repair agent
meets the support agent on the same node.
% -- NARRATE
Both share a 6,144-token prefix of instructions and tool schemas, free of
customer or repository data. For support, that is 75\% of each session and
sharing gives 4$\times$; for repair it is 18.8\% and gives 1.25$\times$. Both
pay a similar tax. The difference is who waits: a 1.3 s rebuild adds 2.6\%
to a 50 s repair cycle that a test harness waits through, while a support
rebuild lands in front of a customer's first token.
% -- WARN
The support inputs are scenario values, not measurements.
% -- FLEX
[CORE]
}

\footnotesize
\begin{tabular}{@{}L{0.34\textwidth}L{0.29\textwidth}L{0.29\textwidth}@{}}
  \toprule
  & \textbf{Repair agent} & \textbf{Support agent} \\
  \midrule
  Retained tokens per session & 32,768 & 8,192 \\
  Shared 6,144-token prefix, fraction & 18.8\% & 75\% \\
  Resident: unshared / shared & 48 / 60 (1.25$\times$) & 195 / 780 (4$\times$) \\
  Generate, then idle wait & 5 s, then 45 s test & 2 s, then 20 s reply \\
  Residency tax if kept & 10$\times$ & 11$\times$ \\
  Resume: reload vs recompute & 21 ms vs 1.3 s & 5 ms vs 0.32 s \\
  Who waits on the resume & the harness, behind a test & a customer, for a first token \\
  \bottomrule
\end{tabular}

\bottomline{Repair: move during tests, drop if host memory is tight (+2.6\% of the cycle). Support: keep or move, never drop.}
\end{frame}

\begin{frame}{Occupied fraction alone cannot tell health from collapse}
\note{
[3 min]
% -- NARRATE
Two pools at the same occupancy: one is mostly dormant prefixes that cheap
reclamation can free, the other is all pinned by active decodes with no room
to grow. Classify every block into exactly one row, in order: a block being
copied out is still pinned; a block allocated for an inbound reload is
reserved. Then watch for churn: blocks evicted and soon reloaded or rebuilt,
while tail latency rises and completed work falls.
% -- LINK
This is thrashing, which time-sharing systems met in the 1960s. Denning's
remedy was load control, not faster paging.
% -- FLEX
[OPTIONAL]
}

\footnotesize
\begin{tabular}{@{}L{0.25\textwidth}L{0.68\textwidth}@{}}
  \toprule
  \textbf{Pool state} & \textbf{Accounting and consequence} \\
  \midrule
  Active or pinned & count each physical block once; not reclaimable while any owner or operation uses it \\
  Reserved or in flight & being populated, copied, or transferred; unavailable for admission \\
  Reclaimable dormant & no active use; can be discarded, at the cost of a future reuse \\
  Genuinely free & no owner, operation, or reservation; allocatable now \\
  \bottomrule
\end{tabular}

\vspace{0.15cm}
\textbf{Churn} (evict, then reload or rebuild soon after) plus rising tail latency under steady load means the working set of active trajectories has overrun memory. The remedy is \textbf{admission control}, not faster eviction.

\bottomline{Admission is judged from block states and queue dynamics, never from a single utilization number.}
\end{frame}

% =============================================================================
% SECTION 8: WRAP-UP
% =============================================================================
\section{Wrap-Up}

\begin{frame}{Fallacies}
\note{
[3 min]
% -- NARRATE
Read each claim, ask for a show of hands on ``true,'' then give the number.
% -- FLEX
[CORE]
}

\small
\textbf{Fallacy:} \textit{Model size, not retained context, decides how many sequences a node can serve.}\\
The same 526.0 GB pool holds 195 sequences at 8,192 tokens and 12 at 131,072, with the model unchanged.

\vspace{0.2cm}
\textbf{Fallacy:} \textit{Compatible cached attention state proves that its evidence is current.}\\
A hit on a 4,096-token prompt saves about 169 ms whether it holds r17 or r18 evidence. Only the harness can tell the difference.

\vspace{0.2cm}
\textbf{Fallacy:} \textit{A resident-fit count is a throughput guarantee.}\\
41 resident slots sustain 0.82 turns/s when held through a 45 s test and 8.2 turns/s when released: residence sets the rate.
\end{frame}

\begin{frame}{Pitfalls}
\note{
[3 min]
% -- NARRATE
Each pitfall is a missing measurement: what is actually used, what is
actually contested, and what is actually the record.
% -- FLEX
[CORE]
}

\small
\textbf{Pitfall:} \textit{Reserving every request's maximum context.}\\
At a 32,768-token cap and a 2,048-token mean, 6.25\% of each reservation holds state; the pool admits 48 instead of about 780 paged snapshots.

\vspace{0.2cm}
\textbf{Pitfall:} \textit{Always moving paused state out of accelerator memory.}\\
Moving r18 costs 42 ms of link transfer and 21 ms on resume. With free blocks in the pool, that buys nothing; under contention it frees 44.98 s of a 45 s wait.

\vspace{0.2cm}
\textbf{Pitfall:} \textit{Treating KV storage as the authoritative task record.}\\
Dropping r18's 10.74 GB costs about 1.3 s of prefill and forfeits no evidence. History, tool outputs, and verification belong in the durable runtime.
\end{frame}

\begin{frame}{What were the key ideas?}
\note{
[3 min] Retrieval practice.
% -- NARRATE
Timer on. Then ask three students for one item each and build the list on
the board. Compare with the next slide.
% -- FLEX
[CORE]
}

\centering
\vspace{0.8cm}
{\Large\bfseries Close your notes.}\\[0.5cm]
{\normalsize Write down \textbf{the bytes-per-token equation}, \textbf{the three reuse gates},\\ and \textbf{the turn rate of 41 slots} with and without a held 45 s wait.}\\[0.3cm]
{\normalsize 90 seconds, no peeking.}
\end{frame}

\begin{frame}{Key takeaways}
\note{
[2 min]
% -- NARRATE
One bullet per section. Return to the subtitle: size and residence time
together set capacity.
% -- FLEX
[CORE]
}

\footnotesize
\begin{itemize}\setlength{\itemsep}{2pt}
  \item \textbf{Footprint:} $\mtok = 2\NL\Hkv\dhead\selem$; fit falls as $1/S$: 195, 48, 12 sequences.
  \item \textbf{Paging:} blocks end reservation slack; copy-on-write forks cost unique blocks only.
  \item \textbf{Prefix cache:} reuse needs compatible, available, permitted state; none checks currency.
  \item \textbf{Prefill:} a hit trades 169 ms for 5 ms; chunking bounds the gap between tokens.
  \item \textbf{Tool waits:} reload beats recompute 62$\times$; wait vs.\ round trip picks the regime.
  \item \textbf{Approximation:} lossy policies change the computation; validate on hard dependencies.
  \item \textbf{Capacity:} $\lam = \Nslot/\Thold$; a held wait costs $1 + \Ttool/\Tmodel$ in turn rate.
\end{itemize}
\vfill\begin{mlsyscard}[top=0.15em,bottom=0.15em,fontupper=\footnotesize]{crimson}\textbf{Today's thesis.} Retained attention state, and how long it stays, sets agent serving capacity.\end{mlsyscard}
\end{frame}

\begin{frame}{The key relations on one page}
\note{
[1 min]
% -- NARRATE
The reference slide students should keep. Each relation answers one
capacity question before anything is deployed.
% -- FLEX
[OPTIONAL]
}

\footnotesize
\begin{tabular}{@{}L{0.30\textwidth}L{0.66\textwidth}@{}}
  \toprule
  \textbf{Quantity} & \textbf{Relation} \\
  \midrule
  Bytes per token, per sequence & $\mtok = 2\NL\Hkv\dhead\selem$, \quad $\Mkv = S\,\mtok$ \\
  KV pool & $\Mpool = M_{\text{HBM}} - M_{\text{weights}} - M_{\text{runtime}}$ \\
  Tail waste per sequence & $0 \le \Delta_{\text{tail}} \le \Sblk - 1$ \\
  Prefix hit rate & $\rho_{\text{prefix}} = \sum S_{\text{match}} / \sum S_{\text{prompt}}$ \\
  Iteration token budget & $\sum C_i + |\mathcal{D}| \le \Nbudget$ \\
  Reload vs recompute & $T_{\text{reload}}/T_{\text{recompute}} = \mtok\Rattain/(2P\,\BWlink)$ \\
  Sharing gain (large pool) & $1/(1-f)$ \\
  Little's law & $\Nslot = \lam\,\Thold$ \\
  Residency tax & $1 + \Ttool/\Tmodel$ \\
  \bottomrule
\end{tabular}
\end{frame}

\begin{frame}{References}
\note{
[0 min] Reference list for further reading.
% -- FLEX
[OPTIONAL]
}

\mlsysrefsmall{Kilburn62}{Kilburn et al., ``One-Level Storage System,'' IRE Trans. Electronic Computers 1962.}
\mlsysrefsmall{Belady66}{Belady, ``A Study of Replacement Algorithms for a Virtual-Storage Computer,'' IBM Sys. J. 1966.}
\mlsysrefsmall{Little61}{Little, ``A Proof for the Queuing Formula: $L = \lambda W$,'' Operations Research 1961.}
\mlsysrefsmall{Pope23}{Pope et al., ``Efficiently Scaling Transformer Inference,'' MLSys 2023.}
\mlsysrefsmall{Kwon23}{Kwon et al., ``Efficient Memory Management for LLM Serving with PagedAttention,'' SOSP 2023.}
\mlsysrefsmall{Zheng24}{Zheng et al., ``SGLang: Efficient Execution of Structured Language Model Programs,'' 2024.}
\mlsysrefsmall{Agrawal24}{Agrawal et al., ``Taming Throughput-Latency Tradeoff ... with Sarathi-Serve,'' OSDI 2024.}
\mlsysrefsmall{Abhyankar24}{Abhyankar et al., ``InferCept: Efficient Intercept Support for Augmented LLM Inference,'' 2024.}
\mlsysrefsmall{Xiao24}{Xiao et al., ``Efficient Streaming Language Models with Attention Sinks,'' ICLR 2024.}
\end{frame}

\begin{frame}{Next lecture: what should outlive the session?}
\note{
[1 min]
% -- LINK
Today the state was recomputable: dropping r18's cache forfeited no
evidence. Next time the question is the information itself.
% -- FLEX
[CORE]
}

\centering
\vspace{0.4cm}
{\large Dropping r18's attention state loses nothing: the tokens can rebuild it.}\\[0.4cm]
{\normalsize The repair's next session still needs the evidence: the r18 candidate, the r17 failure,\\ the passing reproducer, and the integration checks still unresolved.}\\[0.6cm]
\begin{mlsyscard}{crimson}
\centering\normalsize \textbf{Long-Term Memory} asks which information should persist, how it stays current as its sources change, and how a later session retrieves what its task needs.
\end{mlsyscard}
\end{frame}

\end{document}
